\section{The wall associated to a real character}\label{sec:model}

Fix a nontrivial sign character $\chi$ and let $f:Y\to C$ be its connected double cover. Let $E_{\RR}$ be the flat real line on $C$ with holonomy $\chi$ and
let $E=E_{\RR}\otimes_{\RR}\CC$. Its parallel unit frames differ by
signs and give $E$ a conjugation. Set
\[
 N=E\oplus(K_C\otimes E).
\]
Choose a conformal metric $\rho(z)|dz|^2$ on $C$. In a holomorphic
coordinate and a parallel sign frame write a section as $(u,p\,dz)$.
For $-1\le t\le1$ define $D_t$ by subtracting the right-hand
sides of
\begin{equation}\label{eq:character-model}
 \partial_{\bar z}u=\bar p,\qquad
 \partial_{\bar z}p=\rho(z)(t\operatorname{Re}u-i\operatorname{Im}u)
\end{equation}
from the left-hand sides. In the second equation both sides transform
by the squared modulus of the coordinate Jacobian. The first equation
is an equality of $(0,1)$-forms. All entries transform by the frame
sign, so $D_t$ is a global real Cauchy--Riemann operator with fixed
Dolbeault symbol. We write $B_t=f^*D_t$.

\begin{theorem}\label{thm:character}
For every such $f$ there is a path $D_t$, $-1\le t\le1$, with the
following properties.
\begin{enumerate}
\item The operator on $C$ is invertible near zero. At zero the operator
on $Y$ has a transverse real kernel line on which the deck involution
acts by $-1$. Up to real scaling its graph is the only simple
nonzero normal graph in any covering degree.
\item Every $D_t$ with $t>0$ is rigid in all degrees. For each
$M\ge1$ there exists $0<\delta_M\le1$ such that $D_t$ is rigid
through degree $M$ for $-\delta_M\le t<0$. For
$-\delta_M\le t_-<0<t_+\le1$, $d\le M$, and all $k$,
\begin{equation}\label{eq:mainwall}
 L^C_{d,k}(D_{t_+})-L^C_{d,k}(D_{t_-})=
 \begin{cases}
  L^Y_{e,k}(B_1),&d=2e,\\
  0,&d\text{ odd}.
 \end{cases}
\end{equation}
\item There exists $T>0$ such that the parameters at which some
finite unramified cover carries a nonzero kernel are dense in $(-T,0)$.
In particular no open subinterval of $(-T,0)$ is rigid in every degree.
\end{enumerate}
\end{theorem}

\begin{proposition}\label{prop:model}
At $t>0$ the pullback of $D_t$ by every connected stable map of positive
degree is injective. At $t=0$ its kernel on a stable map $a:X\to C$
consists of the parallel real sections of $a^*E_{\RR}$ in the first
summand. It has dimension one precisely when $a$ lifts to $Y$.
The base operator $D_0$ is invertible. Up to real dilation the unique
simple nonzero kernel graph at zero is the embedded double cover $Y$.
The odd kernel of $B_0$ crosses transversely, with positive determinant
sign on the side $t>0$ and negative sign immediately before zero.
\end{proposition}

\begin{proof}
On a nonconstant component of the normalization use a coordinate $w$.
Write $a'=d(z\circ a)/dw$. The first pulled-back equation gives
$\partial_w\bar u=a'p$. Differentiating and using the second gives
\begin{equation}\label{eq:scalar-model}
 \partial_{\bar w}\partial_w\bar u
    =|a'|^2(\rho\circ a)(t\operatorname{Re}u-i\operatorname{Im}u).
\end{equation}
Both identities hold at critical points. For $s=\operatorname{Im}u$
the equation is $\partial_{\bar w}\partial_w s=|a'|^2(\rho\circ a)s$.
Pairing with $s$ and integrating by parts gives $s=0$ for all $t$.
The integral is defined for the flat real line because the square of a
frame sign is one. The real part satisfies the same equation with
the potential multiplied by $t$. It therefore vanishes when $t>0$
and is parallel when $t=0$. The first equation forces $p=0$ off the
finite critical set and hence everywhere in either case. On a
contracted normalization the kernel sections are holomorphic maps to
a fixed complex vector space and are constant. Matching at the nodes
gives the asserted global kernel, including the monodromy condition
around cycles in the dual graph. Conversely, a nonzero parallel
section of $a^*E_{\RR}$ selects a unit sign over every point of $X$,
which is precisely a lift of $a$ to the unit-frame double $Y$.
For any $r>0$ the two values $\pm r$ in a local frame form a single
embedded copy of $Y$ in $N$. Every nonzero kernel graph has this image
for its value of $r$. On $C$ itself $\chi$ is nontrivial, so the
kernel vanishes. Riemann--Roch and $\deg N=2g-2$ give real index
zero.

Trivialize $f^*E_{\RR}$ by its tautological frame. The kernel of
$B_0$ is generated by $\xi=(1,0)$, negated by the deck involution.
The cokernel has real dimension one by Riemann--Roch on $Y$.
Use the isomorphism $f^*K_C=K_Y$ in the second summand. The functional
on the target of $B_0$ obtained by integrating the real part of the
second coefficient against $dx\,dy$ annihilates the derivative term
by Stokes' theorem and the purely imaginary zeroth-order term
pointwise. On $\dot B_0\xi$ its value is minus the area of $Y$ for
the pulled-back metric. Thus the derivative induces an isomorphism
from the kernel to the cokernel, and the determinant sign changes.

At $t=1$ both zeroth-order terms are off-diagonal and conjugate-linear.
Scaling them to zero gives a complex-linear homotopy for the complex
structure $(i,-i)$ on the two summands. At its split Dolbeault end,
conjugating the second kernel and cokernel changes the usual complex
orientation of the real determinant by
$(-1)^{\operatorname{ind}_{\CC}\bar\partial_{K_Y}}=(-1)^{g(Y)-1}=1$.
The invertible $B_1$ thus has positive sign, which persists for
$0<t\le1$ by injectivity.
\end{proof}

\subsection{Isolation in bounded degree}
The isolation in bounded degree follows from a more general property
of an embedded unramified double. For a section $\xi$ of $f^*N$, let
$\Gamma_r$ be the image in $N$ of the graph of $r\xi$ over $Y$.

\begin{proposition}\label{prop:interval}
Let $D_t$ be a smooth path of real Cauchy--Riemann operators with
fixed symbol on a complex bundle $N\to C$ and let $M\ge1$. Suppose
$B_t=f^*D_t$ is invertible for sufficiently small $t\ne0$, and that
$\ker B_0$ contains a section $\xi$ odd under the deck involution
whose graph over $Y$ is embedded in $N$. Assume that for every holomorphic map to $C$ of positive degree at
most $M$ from a smooth connected curve, allowing arbitrary branching,
the image of each
nonzero kernel graph at zero is some $\Gamma_r$ with $r\ne0$.
Then there exists $\delta_M>0$ such that $a^*D_t$ is injective on
every connected stable map $a$ of degree between one and $M$ whenever
$0<|t|\le\delta_M$.
\end{proposition}

\begin{proof}
If a kernel section is nonzero on a nodal domain, constancy and
matching on contracted components imply that it is nonzero on
some positive-degree normalization. Suppose therefore that
$t_n\ne0$ tends to zero and there are smooth connected covers
$a_n:X_n\to C$ of degrees at most $M$ with nonzero kernel sections
$s_n$. After passing to a subsequence their degrees have one value $d$.

The graphs are pseudoholomorphic for the almost complex structures
$J_{t_n}$ of the linear normal equation. In a bundle chart the
projection is holomorphic, the vertical structure is the given
complex structure and the horizontal-to-vertical part is given by the
zeroth-order term of the operator. Fix a Hermitian metric and a
unitary connection on $N$. On a disk bundle there is a common taming
form $\Omega$ obtained by adding the differential of the fiberwise
Liouville form, extended by the connection, to the pullback of an area
form on $C$. Along the zero section all the structures split, so a
fixed small disk bundle suffices for $t$ near zero. By real homogeneity
we may rescale each $s_n$ to have the same positive maximum norm
$\eta$, with the closed $\eta$-disk inside this bundle. The image
cycle, counted with mapping multiplicity, has $\Omega$-area
$d\operatorname{area}(C)$ by Stokes' theorem. Its mass in a smooth
family of Hermitian metrics is therefore bounded.

The compactness theorem for pseudoholomorphic cycles
\cite[Definition~2.6 and Theorem~2.11]{diw} gives a subsequence
converging as currents and in the Hausdorff topology of supports to
a $J_0$-holomorphic cycle. The compact disk containing all supports
allows application of the theorem on the open taming neighborhood.
The limiting support is connected and contains a point of height
$\eta$. Each of its irreducible curves has positive projection degree
since the complex vector fibers contain no compact complex curves.
Integrating the pullback of the normalized area form, cut off outside
the disk containing all supports, shows that the sum of those degrees
with multiplicity is $d$. Each simple parametrization therefore
solves the central equation over a cover of degree at most $M$. By
hypothesis its image is the zero section or some $\Gamma_r$.
Embeddedness and oddness give $\xi(y)\ne0$ and
$\xi(\iota y)=-\xi(y)$, where $\iota$ is the deck involution.
Consequently these images are pairwise disjoint after identifying
$r$ with $-r$. There are only finitely many components. The connected
limit meeting height $\eta$ must be supported on a single
$\Gamma_r$, $r>0$, possibly with multiplicity.

The two values of $r\xi$ in each fiber have distance at least
$2r\min_Y|\xi|>0$. Hausdorff convergence and compactness imply
that $s_n(x)$ is uniformly close to one of the two values
\emph{over $a_n(x)$}: in a local trivialization the two sheets vary
continuously with the base point. Choosing the nearer sheet thus
defines a continuous lift $b_n:X_n\to Y$ with
\[
    f b_n=a_n,\qquad
    \sup_{x\in X_n}|s_n(x)-r\xi(b_n(x))|\longrightarrow0 .
\]
The lift is holomorphic, since $f$ is locally biholomorphic.
At regular values of $b_n$ average $s_n$ over its sheets to obtain
a section $\sigma_n$ of $f^*N$ on $Y$. Every local inverse sheet
solves $B_{t_n}s=0$ there, so their average does also by real
linearity. This average is bounded at every branch value. Integrating
against test forms away from disks of radius $\delta$ about the branch
values gives boundary errors of order $O(\delta)$, so the extension
is a distributional solution and is smooth by elliptic regularity.
The displayed estimate implies
$\|\sigma_n-r\xi\|_\infty\to0$. In particular $\sigma_n$ is a
nonzero element of $\ker B_{t_n}$ for large $n$, contrary to
invertibility.
\end{proof}

\subsection{Singular parameters on cyclic covers}
For the path \eqref{eq:character-model}, all the hypotheses hold for
each $M$ by Proposition~\ref{prop:model}. The interval of isolation
cannot be chosen uniformly in $M$.

\begin{proposition}\label{prop:spectrum}
There exists $0<T\le1$ such that the parameters $t\in[-T,0)$ at
which $D_t$ has a kernel on a connected unramified cover of $C$ are
dense in $[-T,0]$. In particular there are connected unramified
covers of degrees $2n$ with singular parameters $-\lambda_n$, where
$0<\lambda_n\le A/n^2$ for all sufficiently large $n$.
\end{proposition}

\begin{proof}
Represent a primitive class in $H^1(Y;\ZZ)$ by a smooth map
$\theta:Y\to\RR/\ZZ$, and write $\alpha=d\theta$ for the associated
closed real one-form with integral periods. For $0\le s\le1/2$ use
the flat unitary connection $\nabla_s=d-2\pi i s\alpha$ on the trivial
smooth complex line on $Y$. Its holonomy around a loop $\gamma$ is
$\exp(2\pi i s\int_\gamma\alpha)$. Let $\lambda(s)$ be the lowest
eigenvalue of $\tfrac14\nabla_s^*\nabla_s$ for the pullback of the
metric on $C$. This normalization gives
$-\rho^{-1}\partial_{\bar w}\partial_w$ in a parallel frame for a
metric $\rho|dw|^2$. The common form domain is $H^1(Y;\CC)$. By the
Rayleigh principle $2\sqrt{\lambda(s)}$ is the infimum of
$\|\nabla_s u\|_{L^2}$ over the unit $L^2$ sphere in this domain.
The triangle inequality gives
\[
 |\sqrt{\lambda(s)}-\sqrt{\lambda(s')}|
 \le \pi\|\alpha\|_\infty |s-s'| .
\]
Thus $\lambda$ is continuous, with $\lambda(0)=0$. For
$0<s\le1/2$ the holonomy is nontrivial by primitivity. A
zero eigenfunction would be parallel, so $\lambda(s)>0$.

For rational $s=p/n$ in lowest terms take the connected cyclic
$n$-fold cover $Y_n\to Y$ classified by $[\alpha]$ modulo $n$.
There $\theta$ lifts to $\widetilde\theta:Y_n\to\RR/n\ZZ$ and
$\exp(2\pi i s\widetilde\theta)$ is a global parallel frame for the
pulled-back connection. Express a pulled-back lowest eigensection in
this frame and choose a nonzero real or imaginary part $u$. It is a
real scalar eigenfunction on $Y_n$ with eigenvalue $\lambda(s)$.
In the parallel sign frame induced from $Y$, take $u$ as the first
entry of a section over the composite $a:Y_n\to C$, and use the
isomorphism $a^*K_C\cong K_{Y_n}$ to take $\partial u$ as the second
entry. The first equation of \eqref{eq:character-model} holds by
construction. The second follows from \eqref{eq:scalar-model} at
$t=-\lambda(s)$, since $a'$ is everywhere nonzero.

Set $T=\min\{1,\lambda(1/2)\}$. By continuity the values of $\lambda$
on rationals in $(0,1/2]$ are dense in its image on $[0,1/2]$,
which contains $[0,T]$. This gives the asserted density by restricting
to $0<\lambda(s)<T$ and taking closures. Finally the constant trial
section at $s=1/n$ gives
\[
 0<\lambda_n=\lambda(1/n)
 \le \frac{\pi^2\|\alpha\|_{L^2(Y)}^2}
          {n^2\operatorname{area}(Y)} .
\]
The resulting singular parameters lie on covers of $C$ of degree
$2n$, and belong to the parameter interval for large $n$.
\end{proof}

For sufficiently large $n$, the first negative singularity on the
specified cyclic cover has even multiplicity. Denote by $\beta$ the harmonic representative of the
chosen primitive integral class on $Y$.

\begin{corollary}\label{cor:cyclic}
For all sufficiently large $n$ the negative singular parameter closest
to zero on the specified cover $Y_n\to C$ satisfies
\[
 t_n=-\frac{\pi^2\|\beta\|^2_{L^2(Y)}}{\operatorname{area}(Y)}\,n^{-2}
          +O(n^{-4}).
\]
The normal kernel on this map has real dimension two. Its determinant
sign agrees on the two sides of $t_n$ sufficiently close to $t_n$.
\end{corollary}
\begin{proof}
Use a unitary gauge to replace $\alpha$ by $\beta$ in the flat
connection. For completeness we recall the bottom-eigenvalue
calculation for twisted Laplacians
\cite[Remark~2.3 and Lemmas~4.4--4.5]{gi}. On the fixed Sobolev
domain $H^2(Y;\CC)$ put $H(s)=\nabla_s^*\nabla_s/4$.
The simple eigenvalue zero at $s=0$ continues to a simple real
analytic eigenvalue $\lambda(s)$: integration of the resolvent on a
small circle about zero gives an analytic rank-one projection.
Conjugation makes the eigenvalue even. Since $d^*\beta=0$,
$H'(0)1=0$, and the quadratic term in $H(s)1$ is
$\pi^2s^2|\beta|^2$. Differentiating the eigenvalue equation at
zero and pairing with the normalized constant therefore gives
\[
 \lambda(s)=\frac{\pi^2\|\beta\|^2_{L^2(Y)}}{\operatorname{area}(Y)}
                       s^2+O(s^4).
\]
In particular $\lambda$ is strictly increasing immediately to the
right of zero. The remaining spectrum stays away from zero there.
The spectrum on $Y_n$ is the union, with multiplicities, of the
spectra on the cyclic character spaces $s=j/n$, which are the
spectra of $H(j/n)$ on $Y$. The bottom eigenvalue is periodic in
$s$ by integral holonomy, and bounded below away from zero on compact
subintervals of $(0,1)$. Thus for large $n$ its first positive value
on $Y_n$ is $\lambda(1/n)$, attained only by the two conjugate
characters $j=1,n-1$. The real scalar eigenspace has dimension two.

On the unramified map $Y_n\to C$, equation~\eqref{eq:scalar-model}
and the reconstruction by $\partial u$ identify the real normal
kernel at $t_n=-\lambda(1/n)$ with this eigenspace. Its crossing map
to the cokernel is injective: if the derivative applied to a kernel
vector with first entry $u$ lay in the range, the first component
of a preimage would have real part $w$ satisfying
$(\Delta_{Y_n}/4-\lambda(1/n))w=u$, with $\Delta_{Y_n}$ the
nonnegative scalar Laplacian. Self-adjointness gives $\int_{Y_n}u^2=0$.
The normal operator has index zero, hence its crossing map is an
isomorphism of real two-dimensional spaces. The two sign changes
cancel in its determinant.
\end{proof}

\begin{proof}[Completion of the proof of Theorem~\ref{thm:character}]
The central sign line and its pullbacks in Proposition~\ref{prop:model}
satisfy the central kernel hypothesis of the oriented comparison in
Section~\ref{sec:oriented}. Its reversal has multiplier two on the
liftable stratum, whose stack of lifts has degree two. The interval
in Proposition~\ref{prop:interval} permits continuation to the
stated endpoints. The singular parameters on the other side are
given by Proposition~\ref{prop:spectrum}.
\end{proof}
